\chapter{Reading a Mammogram}

\label{ch:read_mammogram}

Mammography is the primary imaging modality for breast cancer screening and diagnosis. Reading a mammogram requires a systematic approach to detect subtle abnormalities, compare symmetric tissue patterns between breasts, and utilize the standardized \textbf{BI-RADS (Breast Imaging-Reporting and Data System)} lexicon.

---

\section*{1. Mammographic Projections}
A standard screening mammogram consists of two views of each breast:
\textbullet \textbf{Craniocaudal (CC) View}: An axial projection taken from superior to inferior. The outer (lateral) aspect of the breast is at the top of the image, and the medial aspect is at the bottom.
\textbullet \textbf{Mediolateral Oblique (MLO) View}: An oblique projection taken at approximately a 45-degree angle. This view is critical as it captures the maximum amount of breast tissue, including the upper outer quadrant, the \textbf{axillary tail (of Spence)}, and the pectoralis major muscle profile.

---

\section*{2. Breast Density (BI-RADS Composition Categories)}
Breast density refers to the relative proportion of radiodense fibroglandular tissue to radiolucent fat. Higher breast density reduces the sensitivity of mammography because dense tissue (white) can mask small cancers (also white):

\textbullet \textbf{Category A (a)}: The breasts are almost entirely fatty. (High sensitivity).
\textbullet \textbf{Category B (b)}: There are scattered areas of fibroglandular density.
\textbullet \textbf{Category C (c)}: The breasts are heterogeneously dense, which may obscure small masses.
\textbullet \textbf{Category D (d)}: The breasts are extremely dense, which lowers the sensitivity of mammography. (Alternative or supplemental screening like ultrasound or MRI may be indicated).

---

\section*{3. Mammographic Findings: What to Look For}
When examining mammograms, compare the left and right breasts side-by-side to identify asymmetries, and look for these primary findings:

\subsection*{Masses}
A mass is a space-occupying lesion seen in two different projections. Evaluate three characteristics:
1. \textbf{Shape}: Round, oval, or irregular (irregular shapes are highly suspicious).
2. \textbf{Margins}: Circumscribed (benign), obscured, microlobulated, ill-defined, or \textbf{spiculated} (lines radiating from a central mass, highly characteristic of invasive breast cancer).
3. \textbf{Density}: High-density masses appear brighter white than the surrounding fibroglandular tissue and are more suspicious.

\subsection*{Calcifications}
Calcifications are tiny calcium deposits in the breast tissue. They are categorized based on their shape and distribution:
\textbullet \textbf{Typically Benign}: Coarse or "popcorn-like" (degenerating fibroadenoma), large rod-like (duct ectasia), round, or rim/eggshell (calcified cysts).
\textbullet \textbf{Suspicious Calcifications}:
  \textbullet *Amorphous*: Small, hazy, indistinct deposits.
  \textbullet *Fine Pleomorphic*: Variable shapes and sizes, indicating necrotic cellular debris.
  \textbullet *Fine Linear or Branching*: Calcium lining the ducts, highly suspicious for \textbf{Ductal Carcinoma In Situ (DCIS)}.

\subsection*{Architectural Distortion}
The normal parenchymal tissue planes are distorted, with radiating spicules or tethering, but no clear mass is seen. This is a subtle but highly significant finding that can indicate invasive lobular carcinoma, radial scar, or post-surgical scarring.

---

\section*{4. BI-RADS Assessment Categories}
Every mammogram report must conclude with a single BI-RADS category to guide clinical action:

\begin{center}
\begin{longtable}{@{}lll@{}}
\toprule
\textbf{Category} & \textbf{Assessment} & \textbf{Management Action} \\
\midrule
\textbf{BI-RADS 0} & Incomplete & Needs additional imaging (e.g., spot compression, ultrasound) \\
\textbf{BI-RADS 1} & Negative & Routine annual screening \\
\textbf{BI-RADS 2} & Benign finding & Routine annual screening (e.g., stable cyst, benign calcification) \\
\textbf{BI-RADS 3} & Probably benign & Short-interval follow-up (typically 6 months; <2\% cancer risk) \\
\textbf{BI-RADS 4} & Suspicious & Tissue biopsy recommended (divided into 4A, 4B, 4C) \\
\textbf{BI-RADS 5} & Highly suggestive of malignancy & Appropriate action (biopsy \& staging; $>95\%$ cancer risk) \\
\textbf{BI-RADS 6} & Known biopsy-proven cancer & Confirm surgical planning or therapeutic response \\
\bottomrule
\end{longtable}
\end{center}

---

\section*{5. Structured Reporting Example: Suspicious Mass (BI-RADS 4)}

\begin{testbox}[Indications:]
62-year-old female presents with a new palpable lump in the upper outer quadrant of the right breast.
\end{testbox}
\begin{refbox}[Findings:]
\textbf{Breast Composition}: The breasts are heterogeneously dense (BI-RADS Category C), which may obscure small masses.\\
\textbf{Right Breast}: In the upper outer quadrant, 4 cm from the nipple, corresponding to the palpable area, there is an irregular, high-density mass measuring 1.8 x 1.4 cm. The mass exhibits spiculated margins. There are no associated microcalcifications. No skin thickening or nipple retraction is present. A normal lymph node is visualized in the right axilla.\\
\textbf{Left Breast}: No masses, suspicious calcifications, or areas of architectural distortion. No axillary lymphadenopathy.
\end{refbox}
\begin{pearlbox}[Impression:]
1. Palpable, irregular, spiculated right breast mass.
2. BI-RADS Category 4: Suspicious. Ultrasound-guided core needle biopsy of the right breast mass is recommended.
\end{pearlbox}

\clearpage
